% ------------------------------------------------------------------------
% Senac Tex: Modelo de Trabalho Academico para o Centro Universitario Senac
% ------------------------------------------------------------------------

% ========================================================================
% CONFIGURAÇÃO DO DOCUMENTO
% ========================================================================

\documentclass[
    12pt,
    openright,
    oneside,
    a4paper,
    english,
    brazil
]{abntex2}

% ---
% Pacotes básicos
% ---
\usepackage{lmodern}
\usepackage[T1]{fontenc}
\usepackage[utf8]{inputenc}
\usepackage{lastpage}
\usepackage{indentfirst}
\usepackage{color}
\usepackage{graphicx}
\usepackage{microtype}
\usepackage{listings}
\usepackage{indentfirst}
\usepackage[top=3cm, bottom=2cm, left=3cm, right=2cm]{geometry}
\usepackage{float}
\usepackage{caption}

% ---
% Pacotes de citações
% ---
\usepackage[brazilian,hyperpageref]{backref}
\usepackage[alf]{abntex2cite}

% Configurações do backref
\renewcommand{\backrefpagesname}{Citado na(s) página(s):\~}
\renewcommand{\backref}{}
\renewcommand*{\backrefalt}[4]{
    \ifcase #1
        Nenhuma citação no texto.
    \or
        Citado na página #2.
    \else
        Citado #1 vezes nas páginas #2.
    \fi
}

% ========================================================================
% DADOS PARA CAPA E FOLHA DE ROSTO
% ========================================================================

\titulo{Modelagem Baseada em Agentes (ABM) para Simular a Disseminação da Dengue}
\autor{Henrique Kioshi Yamauchi \\\ Renato Jorge Alpalhão}
\local{São Paulo - Brasil}
\data{2025}
\orientador{Prof. Afonso Cesar Lelis Brandão}
\instituicao{
    Centro Universitário Senac - Santo Amaro
    \par
    Bacharelado em Ciência da Computação
}
\tipotrabalho{Trabalho de Conclusão de Curso}
\preambulo{
Monografia apresentada na disciplina Trabalho de Conclusão de Curso, como parte dos requisitos para obtenção do título de Bacharel em Ciência da Computação.
}

% ========================================================================
% CONFIGURAÇÕES DO PDF
% ========================================================================

\definecolor{blue}{RGB}{41,5,195}

\makeatletter
\hypersetup{
    pdftitle={\@title},
    pdfauthor={\@author},
    pdfsubject={\imprimirpreambulo},
    pdfcreator={LaTeX with abnTeX2},
    pdfkeywords={dengue}{modelagem baseada em agentes}{GAMA}{epidemiologia}{modelagem espacial},
    colorlinks=true,
    linkcolor=blue,
    citecolor=blue,
    urlcolor=blue,
    bookmarksdepth=4
}
\makeatother

\setlength{\parindent}{1.25cm}
\setlength{\parskip}{0.2cm}
\SingleSpacing

\makeindex

% ========================================================================
% INÍCIO DO DOCUMENTO
% ========================================================================

\begin{document}

\frenchspacing

% ========================================================================
% CAPA
% ========================================================================

\imprimircapa

% ========================================================================
% FOLHA DE ROSTO
% ========================================================================

\imprimirfolhaderosto\*
\cleardoublepage

% ========================================================================
% RESUMO (EM BRANCO PARA PREENCHIMENTO POSTERIOR)
% ========================================================================
\begin{resumo}

\vspace{8cm}

\begin{center}
\textbf{RESUMO}\\[0.5cm]
\textit{O presente Trabalho de Conclusão de Curso apresenta o desenvolvimento e validação de um modelo híbrido baseado em Modelagem Baseada em Agentes (ABM) para análise da disseminação da dengue em ambiente urbano. O modelo integra dados climáticos reais, informações espaciais e técnicas de aprendizado de máquina, permitindo simular interações entre humanos, vetores e ambiente. Os resultados demonstram alta coerência com séries reais observadas, indicando o potencial da abordagem para apoio ao monitoramento epidemiológico.}
\end{center}

\vspace{6cm}

\textbf{Palavras-chave:} Dengue, ABM, Mesa, Epidemiologia Computacional, Modelagem Espacial.

\end{resumo}
\cleardoublepage

% ========================================================================
% ABSTRACT (EM BRANCO PARA PREENCHIMENTO POSTERIOR)
% ========================================================================
\begin{resumo}[Abstract]
    \begin{otherlanguage*}{english}

\vspace{8cm}

\begin{center}
\textbf{ABSTRACT}\\[0.5cm]
\textit{This thesis presents the development and validation of a hybrid Agent-Based Modeling (ABM) approach to simulate the spread of dengue fever in an urban environment. The model integrates real climate data, spatial information, and machine learning techniques to represent interactions between humans, vectors, and the environment. Results show high consistency with real observed epidemiological data, demonstrating the potential of the approach for public health monitoring and decision support.
}
\end{center}

\vspace{6cm}

\textbf{Keywords:} Dengue, ABM, Mesa, Computational Epidemiology, Spatial Modeling.

    \end{otherlanguage*}
\end{resumo}
\cleardoublepage


% ========================================================================
% FICHA CATALOGRÁFICA (EM BRANCO — BIBLIOTECA INSERE)
% ========================================================================
\begin{fichacatalografica}
    \vspace*{\fill}
    \noindent
    \fbox{
        \begin{minipage}[c][8cm][t]{12cm}
            \textbf{FICHA CATALOGRÁFICA}\\[0.5cm]
            (Inserida posteriormente pela biblioteca do Senac)
        \end{minipage}
    }
    \vspace*{\fill}
\end{fichacatalografica}
\cleardoublepage

% ========================================================================
% FOLHA DE APROVAÇÃO (EM BRANCO)
% ========================================================================
\begin{folhadeaprovacao}

\begin{center}
    {\ABNTEXchapterfont\Large\textbf{\imprimirtitulo}}
\end{center}

\vspace{0.5cm}

Trabalho de Conclusão de Curso aprovado em \textbf{\_\_\_ de \_\_\_\_ de \_\_\_\_}.

\vspace{1.5cm}

\begin{center}
Banca Examinadora:
\end{center}

\vspace{1cm}

\assinatura{\textbf{Prof. Afonso Cesar Lelis
Brandão} \\ Orientador \\ Centro Universitário Senac}
\assinatura{\textbf{Prof(a). \_\_\_\_\_\_\_\_\_\_\_\_} \\ Membro da Banca \\ Centro Universitário Senac}
\assinatura{\textbf{Prof(a). \_\_\_\_\_\_\_\_\_\_\_\_} \\ Membro da Banca \\ Centro Universitário Senac}

\vfill

\begin{center}
 São Paulo, 2026.
\end{center}

\end{folhadeaprovacao}
\cleardoublepage

% ========================================================================
% LISTA DE ILUSTRAÇÕES
% ========================================================================
\pdfbookmark[0]{\listfigurename}{lof}
\listoffigures*
\cleardoublepage

% ========================================================================
% LISTA DE ABREVIATURAS E SIGLAS
% ========================================================================

\begin{siglas}
    \item[ABM] Agent-Based Modeling (Modelagem Baseada em Agentes)
    \item[GIS] Geographic Information System (Sistema de Informação Geográfica)
    \item[GAML] GAMA Modeling Language
    \item[SEIR] Suscetível, Exposto, Infectado e Recuperado
    \item[SEI] Suscetível, Exposto e Infectado
    \item[DENV] Vírus da Dengue
    \item[INMET] Instituto Nacional de Meteorologia
    \item[IBGE] Instituto Brasileiro de Geografia e Estatística
    \item[CSV] Comma-Separated Values
    \item[QGIS] Quantum Geographic Information System
    \item[OMS] Organização Mundial da Saúde
    \item[TCC] Trabalho de Conclusão de Curso
\end{siglas}
\cleardoublepage

% ========================================================================
% SUMÁRIO
% ========================================================================

\pdfbookmark[0]{\contentsname}{toc}
\tableofcontents\*
\cleardoublepage

% ========================================================================
% PARTE TEXTUAL
% ========================================================================

\textual

\chapter{INTRODUÇÃO}

A dengue é uma arbovirose causada pelo vírus da dengue (DENV), cuja transmissão ocorre predominantemente pela picada de fêmeas infectadas de mosquitos do gênero \textit{Aedes}. No Brasil, o \textit{Aedes aegypti} constitui o principal vetor urbano. São reconhecidos quatro sorotipos virais, denominados DENV-1, DENV-2, DENV-3 e DENV-4, cuja circulação representa um desafio à vigilância e à prevenção da doença \cite{brasil, who2025}.

A infecção pode permanecer assintomática, provocar manifestações febris ou evoluir para formas graves. A dinâmica de transmissão depende das relações entre hospedeiros humanos, mosquitos vetores e características do ambiente. O crescimento urbano, a disponibilidade de recipientes favoráveis ao desenvolvimento do mosquito, as condições climáticas e a movimentação das pessoas constituem fatores relevantes para compreender a heterogeneidade espacial e temporal da transmissão \cite{who2025}.

Segundo a Organização Mundial da Saúde, aproximadamente metade da população mundial encontra-se em áreas com risco de dengue. Embora existam medidas preventivas e estratégias de imunização para públicos e condições específicos, a prevenção depende fortemente do controle do vetor, da redução dos criadouros, da vigilância e do acesso oportuno à assistência em saúde. A relevância do problema motiva o uso de abordagens computacionais capazes de explorar mecanismos e cenários sob hipóteses explícitas \cite{who2025}.

O município de São Paulo apresenta um contexto territorial complexo para esse tipo de investigação. A concentração populacional, a diversidade de usos do solo e os deslocamentos entre diferentes regiões tornam relevante a investigação de representações espaciais da dinâmica de transmissão da dengue.

A modelagem baseada em agentes --- \textit{Agent-Based Modeling} (ABM) --- permite representar indivíduos heterogêneos com atributos, estados e comportamentos próprios. As interações entre esses agentes e o ambiente podem gerar padrões coletivos que não são especificados diretamente como saídas agregadas. Plataformas como o GAMA oferecem recursos voltados à implementação de modelos espacialmente explícitos e à execução de experimentos computacionais \cite{macal2010, taillandier2019}.

No contexto da dengue, estudos anteriores demonstram a aplicação da modelagem baseada em agentes para representar as interações entre humanos, mosquitos e ambiente. Garcias (2024), por exemplo, desenvolveu um modelo na plataforma GAMA para Portugal continental, enquanto Oliveira et al. (2025) apresentaram uma simulação multiagentes aplicada ao município de Alto Santo, no Ceará. Esses trabalhos demonstram possibilidades para a utilização dessa abordagem, mas diferenças relacionadas à área geográfica, às espécies vetoriais, aos parâmetros e às características da população exigem que cada aplicação seja definida de acordo com seu contexto.

Diante desse cenário, este trabalho propõe o desenvolvimento e a análise de um simulador espacial e estocástico baseado em agentes para representar a dinâmica de transmissão da dengue no município de São Paulo. O modelo é desenvolvido na plataforma GAMA, utilizando a linguagem GAML, e busca representar as interações entre agentes humanos, mosquitos \textit{Aedes aegypti} e componentes ambientais.

\section{Problema de pesquisa}

A construção de um simulador computacional não se limita à implementação de regras capazes de gerar gráficos ou séries temporais. Para que seus resultados possam ser adequadamente interpretados, é necessário que exista correspondência entre o modelo conceitual, as regras implementadas e os procedimentos utilizados para executar os experimentos.

Em modelos baseados em agentes, aspectos como a definição dos estados dos agentes, as regras de interação, a representação espacial, a escala temporal e os processos estocásticos podem influenciar o comportamento da simulação. Dessa forma, a execução do programa e a plausibilidade visual de suas saídas não são suficientes, isoladamente, para demonstrar a consistência do modelo ou a reprodutibilidade dos resultados.

No contexto da representação da transmissão da dengue no município de São Paulo, surge, portanto, o desafio de desenvolver um simulador que permita representar as interações entre humanos, mosquitos e ambiente, ao mesmo tempo em que suas regras e condições de execução possam ser explicitadas e avaliadas.

Diante disso, estabelece-se a seguinte pergunta de pesquisa: \textit{como desenvolver e avaliar um simulador espacial e estocástico baseado em agentes, utilizando a plataforma GAMA, capaz de representar as interações entre humanos, mosquitos \textit{Aedes aegypti} e ambiente na dinâmica de transmissão da dengue no município de São Paulo, considerando a consistência das regras implementadas e a reprodutibilidade dos experimentos?}

\section{Justificativa}

A utilização da dengue como domínio de aplicação oferece um contexto relevante para a investigação de técnicas de modelagem computacional. A dinâmica da doença envolve agentes heterogêneos, interações locais, eventos estocásticos e fatores ambientais, características que podem ser representadas por meio da modelagem baseada em agentes.

Do ponto de vista da Ciência da Computação, o desenvolvimento do simulador envolve desafios relacionados à especificação de entidades e comportamentos, à implementação de regras de interação, à representação de informações espaciais e à execução de experimentos computacionais. Além disso, a documentação das decisões adotadas e a avaliação do comportamento da implementação são importantes para que o modelo possa ser compreendido, reproduzido e aperfeiçoado.

A escolha da plataforma GAMA se justifica por suas capacidades de representar agentes espacialmente explícitos, incorporar dados territoriais e organizar experimentos computacionais. Entretanto, a utilização de uma plataforma especializada não garante, por si só, a adequação das regras implementadas ou a consistência dos resultados, tornando necessária a análise do modelo e de seus procedimentos de execução \cite{taillandier2019}.

Assim, a contribuição deste trabalho consiste no desenvolvimento e na documentação de um simulador baseado em agentes aplicado ao município de São Paulo, bem como na análise de aspectos relacionados à consistência de sua implementação e à reprodução dos experimentos. O trabalho não pretende estabelecer um sistema de previsão epidemiológica validado, mas contribuir para a construção e a avaliação de uma ferramenta computacional destinada à representação da dinâmica de transmissão da dengue.



\section{Objetivo geral}

Desenvolver e avaliar um simulador computacional espacial e estocástico baseado em agentes, utilizando a plataforma GAMA e dados de Sistemas de Informação Geográfica (SIG), para representar a dinâmica de transmissão da dengue no município de São Paulo por meio das interações entre agentes humanos, mosquitos \textit{Aedes aegypti} e componentes ambientais.

\section{Objetivos específicos}

\begin{enumerate}
    \item Investigar os fundamentos da transmissão da dengue e da modelagem baseada em agentes, identificando elementos relevantes para a especificação do simulador;
    
    \item Especificar os agentes, atributos, estados, comportamentos e interações que compõem o modelo;
    
    \item Desenvolver a representação espacial e temporal do ambiente de simulação, utilizando dados geográficos compatíveis com o recorte territorial adotado;
    
    \item Implementar as regras relacionadas às interações entre humanos, mosquitos e componentes ambientais na plataforma GAMA;
    
    \item Documentar as hipóteses, os parâmetros e os procedimentos adotados na construção e execução do modelo;
    
    \item Realizar experimentos computacionais para analisar o comportamento do simulador e a influência de parâmetros e processos estocásticos nos resultados;
    
    \item Identificar limitações da implementação e propor possibilidades de aperfeiçoamento do modelo.
\end{enumerate}

\section{Delimitação do estudo}

O recorte territorial deste estudo corresponde ao município de São Paulo, Estado de São Paulo, Brasil. A representação computacional do território é definida a partir das bases geográficas e dos dados efetivamente utilizados no desenvolvimento do simulador. Dessa forma, a utilização de regiões agregadas ou de geometrias administrativas não implica, por si só, a reprodução completa da distribuição populacional, da mobilidade ou do risco epidemiológico em todos os locais do município.

O modelo concentra-se na representação das interações relacionadas à transmissão da dengue entre humanos, mosquitos \textit{Aedes aegypti} e componentes ambientais, de acordo com as regras e hipóteses estabelecidas para o simulador.

Este trabalho não tem como objetivo representar simultaneamente os quatro sorotipos do vírus da dengue, reproduzir integralmente os mecanismos imunológicos da doença ou produzir previsões epidemiológicas para o município. Da mesma forma, os resultados obtidos pela simulação não devem ser interpretados automaticamente como equivalentes aos registros oficiais de casos sem procedimentos específicos de calibração e validação.

A avaliação do simulador concentra-se nos aspectos definidos no escopo da pesquisa e nas condições experimentais estabelecidas. Dessa forma, eventuais limitações relacionadas aos dados utilizados, à escala de representação e às simplificações adotadas serão consideradas na interpretação dos resultados.
\chapter{FUNDAMENTAÇÃO TEÓRICA E TRABALHOS RELACIONADOS}

Este capítulo apresenta os fundamentos teóricos necessários à compreensão do problema abordado e ao desenvolvimento de um simulador espacial baseado em agentes para representar aspectos da transmissão da dengue. Inicialmente, são discutidos os principais conceitos epidemiológicos relacionados à doença e ao seu vetor. Em seguida, são apresentados conceitos de modelagem e simulação computacional, com ênfase na modelagem baseada em agentes e nos modelos epidemiológicos representados por estados.

Também são abordados a representação espacial por meio de Sistemas de Informação Geográfica (SIG), os conceitos de estocasticidade e experimentação computacional e os recursos da plataforma GAMA e da linguagem GAML. Por fim, são discutidos aspectos relacionados à documentação, verificação e validação de modelos, além de trabalhos anteriores relacionados à aplicação de simulações baseadas em agentes em problemas epidemiológicos.

\section{Fundamentos epidemiológicos da dengue}

A dengue é uma arbovirose causada pelo vírus da dengue e transmitida principalmente por mosquitos do gênero \textit{Aedes}. A compreensão dos mecanismos básicos de transmissão é necessária para definir quais entidades, estados e interações devem ser considerados em uma representação computacional do fenômeno.

\subsection{Doença, agente etiológico e sorotipos}

O vírus da dengue apresenta quatro sorotipos, denominados DENV-1, DENV-2, DENV-3 e DENV-4. A infecção pode ser assintomática ou apresentar manifestações clínicas de diferentes níveis de gravidade. A circulação de diferentes sorotipos constitui um aspecto relevante da dinâmica epidemiológica da doença, especialmente devido às relações entre infecções anteriores, resposta imunológica e risco de formas graves \cite{who_dengue}.

Do ponto de vista da modelagem computacional, a existência de múltiplos sorotipos demonstra que a representação da recuperação de um indivíduo pode envolver simplificações. Um modelo que represente apenas uma dinâmica epidemiológica genérica pode ser adequado ao seu propósito, desde que suas hipóteses sejam explicitadas e suas limitações sejam consideradas na interpretação dos resultados.

\subsection{Vetor e ciclo de desenvolvimento}

No contexto urbano brasileiro, o \textit{Aedes aegypti} possui papel central na transmissão da dengue. As fêmeas adultas participam diretamente da transmissão do vírus, enquanto o desenvolvimento do mosquito envolve estágios anteriores, como ovo, larva e pupa.

O ciclo de desenvolvimento do vetor é influenciado por condições ambientais, incluindo temperatura e disponibilidade de água. Dessa forma, a representação computacional do mosquito pode envolver tanto sua dinâmica epidemiológica quanto processos relacionados ao seu desenvolvimento e à disponibilidade de criadouros.

Entretanto, o nível de detalhamento adotado em um modelo deve ser compatível com seus objetivos, com os dados disponíveis e com o custo computacional da simulação. Nem todos os processos biológicos precisam ser representados individualmente para que o modelo possa investigar uma questão específica. A escolha das abstrações deve, portanto, ser explicitada durante a especificação do modelo.

\subsection{Transmissão humano--vetor}

A transmissão da dengue envolve interações entre humanos e mosquitos vetores. Um mosquito suscetível pode adquirir o vírus ao realizar alimentação sanguínea em um indivíduo infectante. Após um período de incubação extrínseca, o mosquito pode tornar-se capaz de transmitir o vírus. De maneira semelhante, um indivíduo suscetível pode adquirir a infecção após a picada de um mosquito infectante.

Esses processos evidenciam a necessidade de distinguir conceitos como infecção, período de incubação e infectividade. Um indivíduo ou vetor recém-infectado não necessariamente se torna imediatamente capaz de transmitir o vírus. Em modelos computacionais, essa característica pode ser representada por estados intermediários, permitindo distinguir agentes expostos daqueles que já se encontram em condição infectante.

A ocorrência da transmissão também depende das oportunidades de contato entre hospedeiros e vetores. Em modelos espacialmente explícitos, a localização, a mobilidade e as regras de proximidade podem influenciar a probabilidade de interação entre os agentes.

\subsection{Fatores ambientais e medidas de controle}

A dinâmica da dengue é influenciada por fatores biológicos, ambientais e sociais. Temperatura e precipitação podem afetar diferentes etapas do ciclo do vetor, enquanto a disponibilidade de criadouros e a distribuição espacial da população podem modificar as oportunidades de contato entre humanos e mosquitos.

Além disso, medidas de prevenção e controle podem atuar sobre diferentes componentes do sistema. A redução de criadouros, as intervenções destinadas ao controle vetorial e outras estratégias de prevenção não possuem necessariamente o mesmo mecanismo de ação. Dessa forma, quando representadas em um modelo computacional, essas medidas devem ser associadas a regras compatíveis com o efeito que se pretende investigar.

\section{Modelagem e simulação computacional}

A modelagem computacional consiste na construção de uma representação simplificada de um sistema com o objetivo de investigar determinados aspectos de seu funcionamento. A simulação corresponde à execução das regras definidas para o modelo ao longo do tempo, permitindo observar o comportamento produzido pelas relações entre seus componentes.

É importante distinguir o fenômeno real, o modelo conceitual, a implementação computacional e os resultados produzidos pela simulação. Diferenças entre os resultados esperados e os resultados observados podem estar relacionadas às hipóteses adotadas, aos dados utilizados, às regras implementadas ou aos procedimentos experimentais empregados.

\subsection{Abordagens de simulação}

Diferentes abordagens podem ser utilizadas para representar sistemas dinâmicos. A dinâmica de sistemas enfatiza relações agregadas entre variáveis, estoques e fluxos. A simulação a eventos discretos representa alterações de estado desencadeadas por eventos. Já a modelagem baseada em agentes representa entidades individuais e suas interações com outras entidades e com o ambiente \cite{macal_north_2010}.

\{aq deveria ter uma tabela\}

As abordagens não são necessariamente excludentes. Um modelo baseado em agentes pode utilizar estados epidemiológicos semelhantes aos empregados em modelos compartimentais e produzir indicadores agregados para análise.

\subsection{Modelos explicativos e preditivos}

A finalidade de um modelo influencia as abstrações utilizadas, os dados necessários e os critérios adotados para avaliar seus resultados. Modelos explicativos buscam investigar mecanismos e relações capazes de produzir determinados comportamentos, enquanto modelos preditivos têm como objetivo estimar o comportamento futuro ou desconhecido de um sistema sob determinadas condições.

A reprodução de um padrão observado por uma simulação não demonstra, isoladamente, que o modelo possui capacidade preditiva. Da mesma forma, um modelo utilizado para investigar relações entre mecanismos não precisa necessariamente fornecer previsões epidemiológicas para cumprir seu objetivo.

Essa distinção é relevante para a interpretação dos resultados de modelos epidemiológicos. A adequação de um simulador depende de seu propósito, das hipóteses utilizadas, da qualidade dos dados disponíveis e dos procedimentos empregados para sua avaliação.

\section{Modelagem e simulação baseada em agentes}

A modelagem baseada em agentes, ou \textit{Agent-Based Modeling} (ABM), representa sistemas por meio de entidades capazes de possuir atributos, estados e comportamentos próprios. Os agentes interagem entre si e com o ambiente, e essas interações podem produzir padrões coletivos que emergem da dinâmica do sistema \cite{macal_north_2010}.

Essa abordagem é especialmente adequada para sistemas que apresentam heterogeneidade entre entidades e relações dependentes do contexto local. Na representação da dengue, por exemplo, agentes humanos e mosquitos podem possuir localizações e estados epidemiológicos distintos, enquanto o ambiente pode influenciar suas oportunidades de interação.

\subsection{Entidades, atributos e comportamentos}

Um agente pode ser caracterizado por atributos que descrevem suas propriedades e por variáveis de estado que podem ser modificadas durante a execução da simulação. Seus comportamentos correspondem às regras que determinam como o agente interage com outros agentes e com o ambiente.

A definição dessas entidades e regras constitui uma das principais etapas da modelagem baseada em agentes. O comportamento agregado da simulação resulta das interações estabelecidas no nível individual, podendo produzir padrões temporais ou espaciais que não são definidos diretamente como uma única regra global.

Entretanto, a emergência de um padrão na simulação não representa, por si só, uma comprovação de que o modelo reproduz corretamente o sistema real. A interpretação desses padrões depende das hipóteses adotadas e dos procedimentos utilizados para verificar e avaliar o modelo.

\subsection{Interações e vizinhanças}

Em modelos espaciais, as interações entre agentes podem depender de regras de vizinhança ou proximidade. A definição dessas regras determina quais entidades podem interagir em determinado momento e pode influenciar tanto o comportamento do modelo quanto seu custo computacional.

Uma busca exaustiva entre populações de agentes pode exigir um grande número de comparações. Por esse motivo, a representação espacial e as estratégias utilizadas para identificar agentes próximos constituem aspectos relevantes no desenvolvimento de simulações com grande quantidade de entidades.

\section{Modelos epidemiológicos e máquinas de estados}

Modelos epidemiológicos compartimentais representam indivíduos ou populações por meio de estados relacionados à evolução de uma doença. Entre os modelos mais conhecidos encontram-se o SIR, composto pelos estados suscetível, infectado e recuperado, e o SEIR, que acrescenta um estado de exposição ou incubação entre a suscetibilidade e a condição infectante.

Em uma modelagem baseada em agentes, esses estados podem ser representados individualmente. Dessa forma, cada agente possui uma condição epidemiológica própria e realiza transições conforme regras, eventos, temporizadores ou probabilidades definidos pelo modelo.

\subsection{Representação SEIR e SEI}

Para a população humana, uma estrutura SEIR pode representar os estados:

\begin{itemize}
    \item \textbf{Suscetível (S):} indivíduo que pode adquirir a infecção;
    \item \textbf{Exposto (E):} indivíduo infectado que ainda não se encontra em condição infectante;
    \item \textbf{Infectante (I):} indivíduo capaz de participar da transmissão;
    \item \textbf{Recuperado (R):} indivíduo que concluiu o processo epidemiológico representado pelo modelo.
\end{itemize}

Para os mosquitos, uma estrutura SEI pode representar os estados suscetível, exposto e infectante. A utilização desses estados permite distinguir a aquisição da infecção do momento em que o vetor passa a participar efetivamente da transmissão.

\subsection{Máquinas de estados}

A representação computacional de um modelo epidemiológico pode ser compreendida como uma máquina de estados. Nesse caso, é necessário definir quais estados são válidos, quais transições podem ocorrer e quais condições determinam a passagem de um estado para outro.

A especificação dessas transições contribui para a clareza do modelo e para sua posterior verificação. Uma representação conceitual bem definida permite comparar o comportamento esperado com as regras efetivamente implementadas no software.

Além disso, atributos como vacinação ou outras condições individuais podem ser representados separadamente dos estados epidemiológicos. Essa separação permite definir de maneira explícita como tais características influenciam os eventos de infecção ou outros processos representados pela simulação.

\section{Representação espacial e Sistemas de Informação Geográfica}

Os Sistemas de Informação Geográfica permitem organizar e analisar dados associados a localizações no território. Dados vetoriais podem representar pontos, linhas e polígonos, enquanto dados matriciais podem representar superfícies ou variáveis distribuídas espacialmente.

Em uma simulação espacial, os dados geográficos podem desempenhar funções que ultrapassam a visualização do território. As geometrias podem delimitar o ambiente, orientar a distribuição dos agentes, restringir deslocamentos e apoiar cálculos de distância e proximidade.

A integração entre dados geográficos e modelos computacionais exige atenção à compatibilidade entre sistemas de referência de coordenadas, unidades de distância, qualidade das geometrias e níveis de agregação espacial utilizados.

\subsection{Mobilidade e comportamento espacial}

A mobilidade constitui um dos elementos que podem influenciar as oportunidades de interação entre agentes. A escolha de um destino e a definição da trajetória até esse destino representam problemas distintos e podem ser modeladas de diferentes maneiras.

Quando uma simulação utiliza dados agregados, é necessário explicitar quais simplificações são adotadas para representar o movimento individual. Dados sobre deslocamentos entre regiões, por exemplo, não fornecem automaticamente as trajetórias realizadas por cada agente.

Dessa forma, a mobilidade em modelos baseados em agentes deve ser compreendida como uma representação definida a partir dos objetivos do estudo, dos dados disponíveis e das hipóteses adotadas.

\subsection{Integração entre SIG e ABM}

A integração entre Sistemas de Informação Geográfica e modelagem baseada em agentes permite incorporar informações territoriais diretamente ao funcionamento da simulação. Taillandier e Drogoul (2010) discutem diferentes formas de utilização de dados geográficos, incluindo a leitura de arquivos espaciais, o uso das geometrias como elementos que condicionam operações do modelo e a representação de objetos geográficos como agentes.

Na leitura de arquivos espaciais, as geometrias e seus atributos podem ser utilizadas para inicializar ou delimitar o ambiente. Em outras situações, elementos territoriais podem participar diretamente das regras de interação e comportamento. A escolha entre essas possibilidades depende da forma como o espaço é incorporado ao modelo.

Essa integração é particularmente relevante para estudos epidemiológicos que consideram heterogeneidade territorial, distribuição espacial de populações e relações entre localização e interação.

\section{Estocasticidade e experimentação computacional}

Modelos estocásticos incorporam componentes aleatórios em determinados processos. Em uma simulação epidemiológica, a aleatoriedade pode estar presente, por exemplo, na ocorrência de eventos de transmissão, na movimentação dos agentes ou em processos relacionados ao ciclo de vida dos vetores.

A presença de estocasticidade implica que execuções realizadas com os mesmos parâmetros podem produzir resultados diferentes. Por esse motivo, a análise de modelos estocásticos requer procedimentos que considerem a variabilidade entre execuções.

\subsection{Sementes e reprodutibilidade}

Os mecanismos de geração pseudoaleatória utilizam sementes para controlar a sequência de valores produzidos durante uma execução. A utilização de uma mesma configuração e de uma mesma semente permite, em determinadas condições, reproduzir uma sequência de eventos aleatórios.

O controle das sementes contribui para a reprodutibilidade dos experimentos e para a comparação entre diferentes configurações. Entretanto, a reprodução de uma execução também depende da preservação das condições de inicialização, dos parâmetros utilizados e da versão da implementação.

\subsection{Replicações e análise de variabilidade}

Uma única execução de um modelo estocástico pode não representar adequadamente a variabilidade associada aos processos aleatórios. A realização de múltiplas replicações permite observar diferenças entre execuções e calcular medidas descritivas, como médias e dispersões.

A quantidade de replicações e os indicadores utilizados devem ser definidos de acordo com os objetivos do experimento e com o custo computacional das simulações. A comparação entre cenários também deve considerar se as diferenças observadas são consistentes diante da variabilidade entre execuções.

\subsection{Sensibilidade e incerteza}

A análise de sensibilidade investiga como alterações nos parâmetros de um modelo podem influenciar suas saídas. Essa análise pode ser realizada por meio de diferentes estratégias, desde a variação controlada de parâmetros até métodos globais capazes de considerar interações entre múltiplas variáveis.

Os índices de Sobol constituem uma abordagem para analisar a contribuição dos parâmetros para a variância das saídas de um modelo \cite{sobol_2001}. Entretanto, sua aplicação requer um planejamento experimental compatível com o método e um número adequado de execuções.

Em modelos espacialmente explícitos, a incerteza também pode apresentar distribuição heterogênea no território e variar ao longo do tempo. Gamal et al. (2025) discutem uma abordagem para analisar a propagação espaço-temporal da incerteza em modelos baseados em agentes, destacando a importância de considerar tanto a dimensão espacial quanto a evolução temporal dos resultados.

Essas abordagens demonstram que a análise de um modelo estocástico não deve se limitar à observação de uma única execução. A escolha dos procedimentos de experimentação deve, contudo, ser compatível com os objetivos e com os recursos disponíveis para o estudo.

\section{Plataforma GAMA e linguagem GAML}

A plataforma GAMA foi desenvolvida para a construção e execução de modelos espaciais baseados em agentes. A plataforma oferece recursos para a representação de agentes, ambientes geográficos e experimentos computacionais, sendo especialmente voltada a modelos que utilizam informações espaciais \cite{taillandier_etal_2019}.

A linguagem GAML é utilizada para especificar os componentes da simulação. Entre os elementos presentes na linguagem estão estruturas destinadas à definição de variáveis globais, espécies de agentes, procedimentos de inicialização, comportamentos executados durante a simulação e experimentos.

A utilização de uma plataforma especializada facilita a integração entre agentes e informações geográficas, mas não substitui a necessidade de especificar adequadamente as regras do modelo. A qualidade da representação depende das hipóteses adotadas, da implementação e dos procedimentos utilizados para verificar o funcionamento do simulador.

\subsection{Operações espaciais}

Modelos espaciais podem utilizar operações relacionadas ao cálculo de distâncias, à identificação de vizinhos, à manipulação de geometrias e à definição de trajetórias. Esses recursos são relevantes para simulações em que a localização dos agentes influencia suas possibilidades de interação.

Entretanto, a disponibilidade de determinado recurso em uma plataforma não significa que ele seja automaticamente utilizado por todos os modelos desenvolvidos nesse ambiente. A descrição de um simulador específico deve distinguir os recursos disponíveis na plataforma das operações efetivamente incorporadas à sua implementação.

\section{Documentação, verificação e validação}

A construção de um modelo computacional envolve não apenas sua implementação, mas também a documentação das hipóteses e regras utilizadas e a avaliação de seu comportamento.

O protocolo ODD, cujo nome deriva de \textit{Overview, Design concepts and Details}, fornece uma estrutura para documentar modelos baseados em agentes. O protocolo organiza a descrição de aspectos como propósito do modelo, entidades, escalas, processos, condições de inicialização, entradas e submodelos \cite{grimm_etal_2020}.

A utilização de uma estrutura de documentação contribui para tornar explícitas as decisões de modelagem e facilitar a compreensão e a reprodução do trabalho. Entretanto, a documentação, por si só, não comprova que a implementação executa corretamente as regras especificadas.

\subsection{Verificação}

A verificação está relacionada à análise da correspondência entre a implementação e as regras ou especificações definidas para o modelo. Esse processo pode envolver testes de transições de estado, condições de inicialização, regras de interação e indicadores produzidos durante a execução.

A verificação é importante porque uma simulação pode produzir resultados aparentemente plausíveis mesmo quando determinadas regras são implementadas de forma diferente da especificação conceitual.

\subsection{Validação}

A validação está relacionada à adequação do modelo ao propósito para o qual foi desenvolvido e à sua capacidade de representar aspectos relevantes do sistema estudado. Dependendo dos objetivos da pesquisa, esse processo pode envolver a comparação entre resultados simulados e dados observados.

A validação deve ser distinguida da verificação. Um software pode executar corretamente as regras especificadas e, ainda assim, utilizar hipóteses insuficientes para representar adequadamente o fenômeno real. Da mesma forma, uma aproximação visual entre uma saída simulada e uma série observada não comprova isoladamente a correção dos mecanismos internos do modelo.

A calibração também representa um procedimento distinto, relacionado ao ajuste de parâmetros com base em informações observadas.

\subsection{Rastreabilidade e reprodutibilidade}

A rastreabilidade busca estabelecer relações claras entre os elementos conceituais do modelo e sua implementação. Essa relação permite identificar como entidades, estados e regras especificadas são transformados em componentes computacionais.

A reprodutibilidade, por sua vez, depende da documentação das condições de execução, dos parâmetros, das sementes utilizadas, das versões do software e dos procedimentos experimentais. Esses aspectos são particularmente relevantes em modelos estocásticos, nos quais alterações nas condições de execução podem produzir resultados diferentes.

\section{Trabalhos relacionados}

A análise de trabalhos relacionados permite identificar diferentes estratégias utilizadas na aplicação de modelos baseados em agentes a problemas epidemiológicos e espaciais. Os estudos discutidos nesta seção são utilizados como referências para compreender decisões de modelagem, representação espacial e experimentação, sem que seus parâmetros ou resultados sejam automaticamente transferidos para o presente estudo.

\subsection{Garcias (2024)}

Garcias (2024) apresenta um modelo baseado em agentes para representar a transmissão da dengue no contexto de Portugal Continental, utilizando a plataforma GAMA. O estudo considera agentes humanos, mosquitos e componentes territoriais, além de discutir parâmetros, mobilidade e condições ambientais.

Sua principal contribuição para este trabalho está na organização da documentação do modelo e na relação entre representação geográfica, regras de comportamento e experimentação. Entretanto, as diferenças entre a espécie vetorial, a área de estudo e as condições epidemiológicas impedem a transferência automática de parâmetros e resultados.

\subsection{Oliveira et al. (2025)}

Oliveira et al. (2025) apresentam uma simulação multiagentes aplicada à transmissão da dengue no município de Alto Santo, no Ceará, também utilizando a plataforma GAMA. O estudo representa humanos, mosquitos, ovos e criadouros, além de utilizar informações espaciais e realizar experimentos com diferentes configurações de parâmetros.

O trabalho fornece uma referência relevante para a organização de entidades, parâmetros e cenários experimentais. Entretanto, sua área de estudo, suas hipóteses e suas configurações não correspondem diretamente ao contexto do município de São Paulo.

\subsection{Estudos espaciais aplicados a outras doenças}

Gharakhanlou, Mesgari e Hooshangi (2019) e Gharakhanlou, Hooshangi e Helbich (2020) apresentam modelos espaciais baseados em agentes aplicados à transmissão da malária. Embora o sistema epidemiológico seja diferente, esses trabalhos contribuem para a compreensão da integração entre entidades, espaço, fatores ambientais e intervenções.

A principal contribuição desses estudos para o presente trabalho está nos aspectos metodológicos relacionados à representação espacial. As regras epidemiológicas e os parâmetros específicos da malária, entretanto, não podem ser transferidos diretamente para a representação da dengue.

\subsection{Gamal et al. (2025)}

Gamal et al. (2025) apresentam uma abordagem para analisar a propagação espaço-temporal da incerteza em modelos baseados em agentes. O estudo destaca que a variabilidade dos resultados pode apresentar diferenças entre unidades espaciais e ao longo do tempo.

A contribuição do trabalho para esta pesquisa está na discussão sobre a análise de incerteza e heterogeneidade espacial. A abordagem apresentada não constitui um conjunto de regras epidemiológicas para a dengue, mas fornece referências conceituais para a avaliação de modelos espaciais e estocásticos.

\subsection{Síntese dos trabalhos relacionados}

Os trabalhos analisados demonstram diferentes possibilidades de aplicação da modelagem baseada em agentes em contextos epidemiológicos e espaciais. Os estudos relacionados contribuem para a compreensão da documentação de modelos, da representação de entidades, da integração de dados geográficos e da realização de experimentos.

Entretanto, cada modelo está associado a uma combinação específica de objetivos, hipóteses, dados, área de estudo e procedimentos de avaliação. Dessa forma, os resultados e parâmetros apresentados em estudos anteriores não substituem a necessidade de especificação e avaliação próprias do simulador desenvolvido neste trabalho.

\chapter{METODOLOGIA}

Este capítulo apresenta os procedimentos metodológicos adotados para o desenvolvimento e a análise do simulador computacional. São descritos a caracterização da pesquisa, as etapas realizadas, a área de estudo e os dados utilizados, os procedimentos de modelagem e implementação, as estratégias de verificação e o planejamento dos experimentos.

A metodologia busca estabelecer uma separação entre o modelo conceitual, sua implementação computacional e os resultados obtidos durante as execuções. Essa distinção é importante para evitar que a existência de uma funcionalidade no código ou a produção de uma saída pelo simulador seja interpretada, isoladamente, como evidência de validação epidemiológica.

\section{Caracterização da pesquisa}

O presente trabalho caracteriza-se como uma pesquisa aplicada, de natureza computacional e caráter exploratório. O estudo envolve o desenvolvimento de um simulador espacial e estocástico baseado em agentes para representar aspectos da dinâmica de transmissão da dengue no município de São Paulo.

A pesquisa utiliza técnicas de modelagem e simulação computacional para representar entidades, estados, interações e processos temporais relacionados ao sistema estudado. A modelagem baseada em agentes foi adotada por permitir a representação de entidades heterogêneas e suas interações com outras entidades e com o ambiente.

O trabalho também possui caráter exploratório, uma vez que busca investigar o comportamento do modelo sob condições controladas e analisar as possibilidades e limitações decorrentes das hipóteses, dos dados e das regras adotadas.

A avaliação do simulador é tratada em diferentes níveis. Inicialmente, busca-se verificar a consistência entre as regras especificadas e a implementação computacional. Posteriormente, são realizados experimentos para analisar o comportamento produzido pelo modelo. A eventual comparação entre resultados simulados e dados epidemiológicos observados constitui uma etapa distinta e depende da compatibilidade entre indicadores, escalas espaciais e temporais e demais condições necessárias para essa análise.

\section{Etapas da pesquisa}

O desenvolvimento da pesquisa foi organizado em etapas relacionadas entre si, conforme apresentado a seguir:

\begin{enumerate}
    \item realização da revisão bibliográfica sobre dengue, modelagem e simulação computacional, modelagem baseada em agentes e representação espacial;
    
    \item definição do problema de pesquisa, dos objetivos e do escopo do estudo;
    
    \item obtenção e preparação dos dados geográficos e climáticos utilizados pelo simulador;
    
    \item especificação conceitual das entidades, estados, interações e processos que compõem o modelo;
    
    \item desenvolvimento e organização da implementação utilizando a plataforma GAMA e a linguagem GAML;
    
    \item definição de procedimentos de verificação para analisar regras, estados, parâmetros e saídas do simulador;
    
    \item planejamento e realização dos experimentos computacionais;
    
    \item coleta e análise dos indicadores produzidos pelas simulações;
    
    \item identificação das limitações do modelo e dos procedimentos adotados.
\end{enumerate}

As etapas não devem ser compreendidas como um processo estritamente linear. A construção de modelos computacionais pode exigir revisões nas regras, nos parâmetros e na implementação à medida que novas verificações são realizadas.

\section{Área de estudo e dados utilizados}

A área de estudo corresponde ao município de São Paulo, localizado no Estado de São Paulo, Brasil. A escolha do município está relacionada à dimensão territorial e populacional da cidade e à possibilidade de utilizar informações espaciais na representação do ambiente da simulação.

A representação computacional do território não busca reproduzir integralmente todas as características geográficas, demográficas e epidemiológicas do município. O espaço utilizado pelo simulador é definido a partir das bases geográficas selecionadas e das simplificações necessárias para atender aos objetivos do estudo.

\subsection{Dados geográficos}

Os dados geográficos são utilizados para representar o espaço da simulação e orientar a distribuição ou localização dos componentes do modelo. A preparação desses dados inclui a verificação das geometrias utilizadas e de sua compatibilidade com a área de estudo.

A utilização de uma base geográfica em uma simulação não implica, por si só, uma representação detalhada de todos os elementos presentes no território. A resolução espacial adotada depende das informações disponíveis e das decisões de modelagem.

Os procedimentos de preparação dos dados devem preservar a identificação da fonte utilizada e das transformações realizadas antes de sua incorporação ao simulador. Essas informações são necessárias para garantir a rastreabilidade entre os dados de entrada e a representação espacial adotada.

\subsection{Dados climáticos}

Os dados climáticos são utilizados como entradas para representar condições ambientais que podem influenciar processos relacionados ao vetor e ao ambiente da simulação.

Antes de sua utilização, os dados devem ser organizados em formato compatível com o processamento realizado pelo modelo. Esse procedimento pode incluir a seleção das variáveis de interesse, a padronização das datas e a verificação da consistência dos valores utilizados.

A utilização de dados climáticos históricos não significa que todas as relações entre clima e dinâmica epidemiológica sejam reproduzidas automaticamente. Os efeitos dessas variáveis dependem das regras implementadas e dos parâmetros definidos no modelo.

\section{Procedimento de modelagem}

O procedimento de modelagem consiste na definição da representação conceitual do sistema antes de sua descrição detalhada em código. Essa etapa estabelece quais entidades serão consideradas, quais atributos e estados podem possuir e quais regras representam suas interações.

A especificação do modelo é orientada pelos conceitos discutidos na fundamentação teórica e pelas informações disponíveis para o estudo. As entidades representam componentes considerados relevantes para o objetivo da simulação, enquanto os comportamentos representam abstrações dos processos que se pretende investigar.

A documentação do modelo utiliza os princípios do protocolo ODD como referência para a organização das informações. Dessa forma, são considerados aspectos como o propósito do modelo, as entidades, as escalas temporais e espaciais, os processos, as condições de inicialização, as entradas e os submodelos \cite{grimm_etal_2020}.

A modelagem também envolve a definição explícita das hipóteses e simplificações adotadas. Processos que não são representados diretamente ou que são descritos de forma agregada devem ser identificados para que os resultados sejam interpretados de acordo com o escopo do modelo.

A descrição detalhada das entidades, dos estados, das interações e das regras conceituais é apresentada no Capítulo 4.

\section{Procedimentos de implementação}

A implementação do simulador é realizada utilizando a plataforma GAMA e a linguagem GAML. O desenvolvimento é organizado em componentes destinados à representação dos agentes e dos processos globais da simulação.

Os procedimentos de implementação incluem a organização dos arquivos do projeto, a definição das espécies de agentes, a inicialização do ambiente e das populações e a implementação das regras responsáveis pela evolução da simulação.

Também são implementados mecanismos para o processamento das entradas geográficas e climáticas, o controle da escala temporal, a configuração de parâmetros e o registro das informações produzidas durante as execuções.

A implementação deve manter correspondência com a especificação conceitual do modelo. Quando uma decisão técnica exigir uma adaptação da representação inicialmente definida, essa alteração deve ser documentada para preservar a rastreabilidade entre o modelo e o código.

A arquitetura do projeto e os procedimentos técnicos utilizados na implementação são apresentados detalhadamente no Capítulo 5.

\section{Procedimentos de verificação}

A verificação tem como objetivo analisar se a implementação apresenta comportamento consistente com as regras e condições definidas para o modelo. Esse procedimento não corresponde, isoladamente, à validação epidemiológica dos resultados.

Os procedimentos de verificação são organizados em diferentes aspectos do simulador, incluindo os estados dos agentes, as transições epidemiológicas, as condições de inicialização, o controle temporal, os parâmetros e os indicadores produzidos.

\subsection{Verificação dos estados e transições}

As transições entre estados devem ser analisadas de acordo com a especificação do modelo. Para isso, devem ser verificadas as condições necessárias para cada transição e a impossibilidade de ocorrência de estados incompatíveis com as regras definidas.

No caso dos indicadores epidemiológicos, também devem ser observadas relações de consistência entre as contagens registradas. Quando os estados forem mutuamente exclusivos e representarem a totalidade de uma população, as somas dos indicadores devem corresponder à quantidade de agentes considerada.

\subsection{Verificação das condições iniciais e dos parâmetros}

As condições de inicialização devem ser verificadas para garantir que as populações e os componentes do ambiente sejam criados de acordo com as configurações definidas.

Os parâmetros utilizados durante a execução também devem ser analisados quanto à sua interpretação e aos limites estabelecidos. Alterações em um parâmetro devem produzir efeitos compatíveis com a regra à qual ele está associado, respeitadas as características estocásticas do modelo.

\subsection{Verificação temporal}

A escala temporal utilizada pela simulação deve apresentar correspondência clara entre os ciclos computacionais e os períodos representados no modelo.

A verificação temporal inclui a análise da frequência de atualização dos processos, da leitura dos dados ambientais e do intervalo utilizado para o registro das saídas. Essa etapa busca evitar interpretações incorretas decorrentes de diferenças entre o tempo computacional e o período representado.

\subsection{Verificação das saídas}

Os indicadores produzidos pelo simulador devem ser analisados quanto à consistência de seus valores e à interpretação das informações registradas.

Entre as verificações previstas estão:

\begin{itemize}
    \item ausência de contagens negativas;
    \item consistência entre o total de agentes e os indicadores de estado;
    \item identificação da escala temporal associada aos registros;
    \item correspondência entre os parâmetros utilizados e os resultados registrados;
    \item análise da possibilidade de reprodução das execuções sob condições equivalentes.
\end{itemize}

Os resultados dos procedimentos de verificação são apresentados posteriormente no capítulo de resultados e avaliação.

\section{Protocolo experimental}

Os experimentos computacionais são utilizados para analisar o comportamento do modelo sob condições previamente definidas. Para garantir a interpretação dos resultados, cada experimento deve registrar as configurações utilizadas, os parâmetros alterados e as condições de execução.

O protocolo experimental considera a definição de um cenário de referência e, quando aplicável, a comparação com outros cenários obtidos pela alteração controlada de parâmetros ou condições do modelo.

\subsection{Cenário de referência}

O cenário de referência corresponde à configuração utilizada como base para a análise do comportamento do simulador. Seus parâmetros e condições iniciais devem ser documentados para permitir a interpretação e a reprodução das execuções.

O cenário de referência não deve ser interpretado automaticamente como uma reprodução exata da situação epidemiológica observada no município. Seu propósito é fornecer uma configuração definida para a realização dos experimentos.

\subsection{Sementes e replicações}

Como o simulador possui componentes estocásticos, diferentes execuções podem produzir resultados distintos mesmo quando utilizam os mesmos parâmetros.

Para analisar essa variabilidade, os experimentos devem considerar múltiplas replicações. O controle das sementes pseudoaleatórias permite identificar as condições utilizadas em cada execução e contribui para a reprodutibilidade dos resultados.

A utilização de uma mesma semente, juntamente com parâmetros e condições equivalentes, deve ser utilizada como procedimento de verificação da reprodutibilidade, quando tecnicamente aplicável.

\subsection{Comparação entre cenários}

A comparação entre cenários deve ser realizada mediante alterações explícitas nas configurações do experimento. Os parâmetros modificados e os indicadores utilizados na comparação devem ser registrados para evitar que diferenças nos resultados sejam atribuídas a fatores não documentados.

A escolha dos cenários deve estar relacionada aos objetivos da pesquisa e às hipóteses que se pretende investigar. Caso sejam realizadas análises de sensibilidade, os parâmetros selecionados e os intervalos avaliados devem ser apresentados de forma explícita.

\section{Indicadores e análise dos resultados}

A análise dos resultados é realizada a partir dos indicadores produzidos durante as execuções do simulador. Os dados registrados devem permitir a identificação das condições da simulação e do comportamento das populações ao longo do tempo.

Entre os indicadores de interesse podem estar as quantidades de agentes em diferentes estados epidemiológicos, informações relacionadas às populações vetoriais e variáveis ambientais utilizadas durante a execução.

Os indicadores são analisados de acordo com os objetivos de cada experimento. Para modelos estocásticos, a interpretação não deve se limitar a uma única execução. Sempre que forem realizadas múltiplas replicações, a análise deve considerar a variabilidade observada entre os resultados.

A apresentação dos resultados pode utilizar tabelas, gráficos e medidas descritivas adequadas aos indicadores disponíveis. A análise deve distinguir resultados efetivamente obtidos durante as execuções de comportamentos esperados ou apenas previstos pela especificação do modelo.

As contagens produzidas pelo simulador também não devem ser consideradas automaticamente equivalentes a indicadores epidemiológicos oficiais. Uma comparação com dados observados exige a análise da definição de cada indicador, da cobertura espacial e da compatibilidade temporal entre as fontes.

\section{Limitações metodológicas}

A metodologia adotada apresenta limitações decorrentes das simplificações necessárias para a construção do modelo, da disponibilidade dos dados e das condições de experimentação.

A utilização de uma população sintética, a agregação espacial, a representação simplificada de processos biológicos e epidemiológicos e a definição de parâmetros constituem exemplos de decisões que podem influenciar o comportamento da simulação.

Além disso, a qualidade das conclusões obtidas depende das verificações e dos experimentos efetivamente realizados. A análise da estrutura do código ou a inspeção de resultados isolados não substitui a execução de testes sistemáticos.

O modelo também não deve ser interpretado como um sistema de previsão epidemiológica validado sem procedimentos específicos de calibração e comparação com dados observados. Os resultados devem ser analisados dentro dos limites definidos pelo propósito da pesquisa e pelas hipóteses adotadas.

A explicitação dessas limitações busca evitar extrapolações que ultrapassem o escopo do estudo e fornecer uma base para futuras melhorias na representação do modelo, nos dados utilizados e nos procedimentos de avaliação.



\chapter{MODELAGEM COMPUTACIONAL DO SIMULADOR}

Este capítulo apresenta a modelagem computacional adotada para representar aspectos da dinâmica de transmissão da dengue no município de São Paulo. A descrição concentra-se no modelo conceitual, apresentando o propósito da simulação, a representação espacial, as entidades consideradas, os estados epidemiológicos e as principais interações entre humanos, mosquitos e componentes ambientais.

A modelagem busca estabelecer uma representação simplificada do sistema estudado. Dessa forma, nem todos os processos biológicos, epidemiológicos e territoriais são reproduzidos individualmente. As regras apresentadas representam abstrações definidas de acordo com os objetivos do estudo e com as possibilidades de implementação do simulador.

A implementação técnica dessas regras, a organização dos arquivos do projeto e os mecanismos específicos utilizados na plataforma GAMA são apresentados no capítulo seguinte.

\section{Propósito e escopo do modelo}

O modelo tem como propósito representar, por meio de uma simulação espacial e estocástica baseada em agentes, aspectos das interações entre a população humana, os mosquitos vetores e o ambiente relacionados à transmissão da dengue.

A utilização de agentes permite representar entidades individuais ou componentes do ambiente com atributos e comportamentos próprios. As interações entre essas entidades ocorrem ao longo do tempo e podem produzir alterações nos estados epidemiológicos das populações representadas.

O modelo não tem como objetivo reproduzir integralmente a dinâmica epidemiológica da dengue no município de São Paulo, nem atuar como um sistema de previsão de casos. Os resultados devem ser interpretados como produtos das regras, parâmetros, condições iniciais e simplificações estabelecidas na simulação.

O recorte territorial corresponde ao município de São Paulo. Entretanto, a utilização de uma representação espacial do município não implica a reprodução individual de todos os seus habitantes, imóveis, vias, criadouros ou demais elementos urbanos.

\section{Representação espacial}

O espaço constitui um componente operacional da simulação, pois influencia a localização dos agentes e suas possibilidades de interação. A representação territorial é construída a partir das bases geográficas selecionadas para o projeto.

A utilização de dados geográficos permite delimitar o ambiente e associar informações espaciais aos componentes da simulação. Entretanto, a resolução espacial adotada representa uma simplificação do território real e deve ser interpretada de acordo com os dados efetivamente incorporados ao modelo.

A distribuição dos agentes no ambiente busca estabelecer uma relação entre a população representada e o espaço da simulação. Dependendo das regras adotadas, essa distribuição pode considerar regiões agregadas ou outras unidades espaciais presentes nas bases utilizadas.

A representação espacial também estabelece o contexto para a mobilidade e para as interações de proximidade. Assim, a localização de um agente pode influenciar quais outras entidades podem ser consideradas em determinado evento da simulação.

\section{Entidades do modelo}

O modelo é composto por entidades que representam componentes considerados relevantes para a dinâmica estudada. Entre elas estão agentes humanos, mosquitos vetores e componentes relacionados ao ambiente e à reprodução vetorial.

A separação entre essas entidades permite atribuir características e comportamentos específicos a cada componente do sistema.

\subsection{Pessoas}

Os agentes humanos representam uma população sintética utilizada na simulação. Cada pessoa pode possuir atributos relacionados à sua localização, condição epidemiológica e demais características necessárias às regras do modelo.

A população sintética não corresponde a uma reprodução individual da população real do município. Seu objetivo é fornecer entidades capazes de participar das interações representadas pela simulação.

Os agentes humanos podem alterar sua localização conforme as regras de mobilidade e participar de eventos de transmissão quando as condições definidas pelo modelo são atendidas.

\subsection{Mosquitos}

Os mosquitos representam os vetores envolvidos na transmissão da dengue. Os agentes podem possuir atributos relacionados à sua localização, condição epidemiológica e estágio de desenvolvimento, de acordo com o nível de detalhamento adotado pelo modelo.

Os mosquitos adultos participam diretamente das interações de transmissão representadas na simulação. Os processos relacionados à reprodução e ao desenvolvimento vetorial são tratados por regras específicas associadas aos componentes ambientais e às formas imaturas.

\subsection{Criadouros e componentes ambientais}

Os componentes ambientais representam condições relacionadas à disponibilidade de locais adequados ao desenvolvimento do vetor. Criadouros ou estruturas equivalentes podem armazenar informações utilizadas pelas regras de reprodução e desenvolvimento dos mosquitos.

A representação desses componentes não busca reproduzir individualmente todos os recipientes existentes no município. Eles constituem abstrações utilizadas para representar condições ambientais relevantes para a dinâmica vetorial.

\subsection{Formas imaturas}

A representação das formas imaturas permite incorporar etapas do ciclo de desenvolvimento do mosquito antes da fase adulta. Dependendo das regras adotadas, ovos e outros estágios podem possuir comportamentos ou tempos de desenvolvimento associados às condições ambientais.

A utilização desses componentes permite estabelecer uma relação entre a dinâmica dos criadouros, as condições climáticas e a população de mosquitos adultos. O grau de detalhamento dessas relações depende das abstrações definidas para o simulador.

\section{Estados epidemiológicos}

A representação epidemiológica utiliza estados para descrever a condição dos agentes ao longo da simulação. A definição explícita desses estados permite organizar as regras de transmissão e estabelecer condições para as transições.

\subsection{Estados da população humana}

A população humana é representada conceitualmente por uma estrutura baseada no modelo SEIR, composta pelos estados suscetível, exposto, infectado e recuperado.

O estado suscetível representa indivíduos que ainda podem adquirir a infecção. Após um evento de transmissão, o agente pode passar ao estado exposto, representando o período entre a infecção e a condição infectante.

Após o período definido pelas regras do modelo, o agente pode alcançar o estado infectado e participar dos eventos de transmissão representados pela simulação. Posteriormente, o indivíduo pode alcançar o estado recuperado.

As transições entre esses estados dependem das condições, probabilidades e intervalos temporais definidos pelo modelo.

\subsection{Estados da população vetorial}

A condição epidemiológica dos mosquitos é representada conceitualmente por uma estrutura SEI, composta pelos estados suscetível, exposto e infectado.

Um mosquito suscetível pode adquirir a infecção ao interagir com um humano infectante. Após o período definido para a evolução da infecção, o mosquito pode passar a uma condição infectante e participar da transmissão para agentes humanos suscetíveis.

A representação por estados permite diferenciar o momento da aquisição da infecção do momento em que o vetor passa a participar da transmissão.

\section{Interações e transmissão}

A transmissão da dengue é representada por interações entre agentes humanos e mosquitos. O modelo não considera a transmissão direta entre pessoas como mecanismo de propagação.

Os eventos de transmissão dependem da ocorrência de interações compatíveis com as regras espaciais e epidemiológicas da simulação. Quando um mosquito e um agente humano se encontram em condições epidemiológicas adequadas, pode ocorrer uma tentativa de transmissão.

A ocorrência efetiva do evento pode depender de um componente probabilístico. Dessa forma, a existência de uma oportunidade de contato não implica necessariamente que a transmissão ocorrerá.

A probabilidade utilizada em cada processo constitui um parâmetro do modelo e deve ser interpretada de acordo com a escala temporal e com as condições em que é aplicada.

A representação da transmissão busca reproduzir a estrutura geral da interação humano--vetor. Entretanto, não representa todos os mecanismos biológicos envolvidos em uma picada ou em uma infecção.

\section{Mobilidade e comportamento espacial}

A mobilidade humana é utilizada para modificar a localização dos agentes e, consequentemente, suas possibilidades de interação com outras entidades.

As regras de deslocamento representam uma simplificação da mobilidade real. A definição de horários, destinos, velocidades e trajetórias deve ser interpretada como uma hipótese de modelagem, especialmente quando não são utilizados dados individuais de deslocamento.

A mobilidade pode alterar a distribuição espacial da população ao longo do tempo e modificar as oportunidades de contato entre humanos e mosquitos.

Os mosquitos também podem possuir regras de deslocamento próprias, compatíveis com a escala espacial adotada. Essas regras contribuem para a dinâmica local da população vetorial e para a distribuição das oportunidades de transmissão.

\section{Ambiente, clima e ciclo vetorial}

As condições ambientais constituem um componente da dinâmica do modelo, especialmente nos processos relacionados à disponibilidade de criadouros e ao desenvolvimento dos mosquitos.

Dados climáticos podem ser utilizados para representar variáveis como temperatura e precipitação. Essas informações são incorporadas às regras ambientais de acordo com os procedimentos definidos pelo modelo.

A chuva pode influenciar a disponibilidade ou as condições dos criadouros, enquanto a temperatura pode estar associada a processos relacionados ao desenvolvimento dos mosquitos. A existência dessas relações na simulação não significa que o modelo reproduza integralmente todas as relações observadas entre clima, vetor e transmissão da dengue.

O ciclo vetorial é representado por processos relacionados à reprodução, ao desenvolvimento das formas imaturas e à formação da população adulta. Os tempos e probabilidades utilizados nesses processos constituem parâmetros ou regras de modelagem e devem ser documentados juntamente com suas fontes ou hipóteses.

\section{Tempo e estocasticidade}

A simulação utiliza uma escala temporal composta por ciclos computacionais. A relação entre os ciclos e os períodos representados deve ser explicitada para permitir a interpretação dos eventos e dos indicadores produzidos.

Diferentes processos podem possuir frequências distintas de atualização. Eventos relacionados à mobilidade, à evolução dos estados epidemiológicos, ao ambiente e à coleta dos resultados podem ocorrer em intervalos específicos definidos pelas regras do simulador.

O modelo também possui componentes estocásticos. Eventos como transmissão, movimentação, desenvolvimento ou outras alterações de estado podem utilizar mecanismos de geração pseudoaleatória.

Como consequência, execuções realizadas com os mesmos parâmetros podem produzir resultados diferentes. O controle das sementes e a realização de múltiplas replicações constituem, portanto, procedimentos importantes para a análise do comportamento do modelo.

\section{Inicialização e parâmetros}

A inicialização estabelece as condições existentes no início de uma execução. Esse processo inclui a criação das populações de agentes e a definição dos valores iniciais das variáveis necessárias à simulação.

As condições iniciais influenciam o comportamento posterior do modelo. Por esse motivo, a quantidade de agentes, a distribuição dos estados epidemiológicos e a configuração do ambiente devem ser registradas como parte da descrição de cada cenário experimental.

Os parâmetros permitem alterar aspectos específicos da simulação sem modificar diretamente sua estrutura conceitual. Entre os parâmetros podem estar probabilidades associadas à transmissão, condições relacionadas ao ambiente, configurações temporais e características das populações.

A escolha dos valores dos parâmetros pode ter origem em referências científicas, dados disponíveis ou hipóteses de modelagem. Essas diferentes origens devem ser distinguidas para evitar que uma hipótese operacional seja interpretada como uma estimativa observacional.

\section{Entradas e saídas}

As entradas do simulador correspondem às informações e configurações utilizadas durante a inicialização e a execução. Entre elas podem estar os dados geográficos, os dados climáticos, as condições iniciais e os parâmetros definidos para cada experimento.

As saídas correspondem aos indicadores registrados durante ou após a execução. Esses indicadores podem incluir informações sobre as populações humanas e vetoriais, seus estados epidemiológicos e variáveis ambientais utilizadas pelo modelo.

O registro das saídas permite acompanhar o comportamento temporal da simulação e fornecer dados para os procedimentos de verificação e análise experimental.

Entretanto, a produção de um indicador pelo simulador não significa automaticamente que ele seja equivalente a uma medida epidemiológica observada. A comparação com dados externos exige compatibilidade entre as definições, escalas e períodos considerados.

\section{Hipóteses e simplificações}

A construção do modelo exige a adoção de simplificações. A representação de uma cidade, de sua população e da dinâmica da dengue envolve processos que não podem ser reproduzidos integralmente em uma única simulação.

Entre as principais limitações do escopo do modelo estão a utilização de uma população sintética, a simplificação da mobilidade, a representação parcial das características ambientais e a abstração de processos biológicos e epidemiológicos.

O modelo também não busca representar simultaneamente todos os sorotipos da dengue, a totalidade das respostas imunológicas ou todos os fatores que influenciam a ocorrência da doença.

As regras ambientais e epidemiológicas devem ser interpretadas como hipóteses de representação. Seus resultados dependem da combinação entre a estrutura do modelo, os parâmetros utilizados, as condições iniciais e os componentes estocásticos.

A explicitação dessas simplificações é necessária para estabelecer os limites de interpretação dos resultados e para identificar possibilidades de aperfeiçoamento em trabalhos futuros.


\chapter{ARQUITETURA E IMPLEMENTAÇÃO}

Este capítulo apresenta a arquitetura e a implementação do simulador desenvolvido na plataforma GAMA utilizando a linguagem GAML. Enquanto o capítulo anterior descreveu conceitualmente os componentes do modelo, esta seção apresenta como essas entidades, regras e processos foram organizados na estrutura computacional do projeto.

A implementação é dividida em módulos destinados à organização dos agentes e dos processos globais da simulação. Essa separação busca facilitar a identificação das responsabilidades de cada componente e estabelecer uma relação entre o modelo conceitual e sua implementação.

\section{Arquitetura geral}

O projeto é organizado em arquivos GAML com responsabilidades distintas. O arquivo \texttt{Main.gaml} atua como ponto de entrada do experimento gráfico e reúne elementos relacionados à execução e à visualização da simulação.

Os demais componentes são responsáveis pela definição das entidades e dos processos que compõem o modelo. Entre os arquivos identificados estão os módulos destinados aos agentes humanos, aos mosquitos, ao ambiente, aos criadouros e às formas imaturas, além de um componente global responsável por processos compartilhados.

A organização geral do projeto pode ser representada da seguinte forma:

\begin{itemize}
    \item \texttt{Main.gaml}: ponto de entrada e configuração do experimento;
    \item \texttt{globals.gaml}: variáveis globais, inicialização, controle temporal e processos compartilhados;
    \item \texttt{Pessoa.gaml}: definição dos agentes humanos e seus comportamentos;
    \item \texttt{Mosquito.gaml}: definição dos agentes vetoriais e suas regras;
    \item \texttt{Ambiente.gaml}: representação dos componentes territoriais;
    \item \texttt{Poca.gaml}: regras relacionadas aos criadouros;
    \item \texttt{Imaturo.gaml}: representação das formas anteriores à fase adulta do mosquito.
\end{itemize}

A divisão em módulos permite separar responsabilidades e reduzir a concentração de regras em um único arquivo. Entretanto, os componentes da simulação permanecem interdependentes, pois processos epidemiológicos, ambientais e espaciais compartilham informações globais e interagem durante a execução.

\section{Estrutura dos agentes}

Os agentes representam os principais componentes individuais da simulação. Cada espécie possui atributos e comportamentos próprios, definidos de acordo com sua função no modelo.

A implementação utiliza espécies GAML para representar agentes humanos, mosquitos e outros componentes necessários à dinâmica da simulação. Essa estrutura permite que diferentes entidades possuam regras específicas e interajam dentro do mesmo ambiente espacial.

\subsection{Agentes humanos}

Os agentes humanos são definidos no módulo destinado à espécie \texttt{Pessoa}. Esse componente reúne os atributos e comportamentos relacionados à população humana, incluindo sua condição epidemiológica e suas regras de movimentação.

A condição epidemiológica é utilizada pelas regras que representam a evolução da infecção e as interações com os mosquitos. A separação desses atributos permite que cada agente mantenha seu próprio estado durante a execução.

O módulo também concentra comportamentos relacionados à mobilidade. Dessa forma, alterações na posição dos agentes são tratadas como parte do comportamento individual da população humana.

\subsection{Agentes vetoriais}

Os mosquitos são definidos em um módulo específico, separado da representação da população humana. Essa organização permite atribuir regras próprias aos vetores, incluindo deslocamento, eventos de alimentação, transmissão e processos relacionados à mortalidade.

A condição epidemiológica dos mosquitos é utilizada para determinar sua participação nos eventos de transmissão. Os estados vetoriais devem ser mantidos de forma consistente durante a execução, pois as contagens produzidas pelo simulador dependem da classificação correta de cada agente.

\subsection{Componentes ambientais e formas imaturas}

O ambiente é representado por componentes específicos destinados à utilização das informações territoriais. Os criadouros são tratados em um módulo separado, no qual são concentradas regras relacionadas às condições ambientais.

A representação das formas imaturas é organizada em um componente próprio, permitindo distinguir os processos anteriores à formação dos mosquitos adultos da dinâmica dos agentes vetoriais adultos.

Essa divisão busca manter separadas as responsabilidades relacionadas ao território, aos criadouros e às diferentes etapas do ciclo vetorial.

\section{Inicialização do ambiente e das populações}

A inicialização estabelece as condições necessárias para o início de uma execução. Esse processo inclui o carregamento das informações espaciais, a criação dos componentes ambientais e a formação das populações utilizadas pelo simulador.

O ambiente geográfico é definido a partir de uma base espacial utilizada pelo projeto. As geometrias carregadas fornecem referências para a delimitação do espaço e para a associação entre os agentes e a representação territorial.

Após a configuração do ambiente, são criados os agentes humanos e vetoriais de acordo com as condições iniciais definidas pelo modelo. A quantidade de agentes utilizada representa uma população sintética e não corresponde diretamente à população total do município de São Paulo.

A distribuição inicial dos agentes constitui uma condição importante para os experimentos, pois influencia as interações que podem ocorrer posteriormente. Por esse motivo, os parâmetros e condições utilizados na inicialização devem ser registrados juntamente com as configurações das simulações.

\section{Implementação da dinâmica epidemiológica}

A dinâmica epidemiológica é implementada por meio de atributos de estado e regras responsáveis pelas transições dos agentes humanos e vetoriais.

Nos agentes humanos, a condição epidemiológica é utilizada para representar as diferentes etapas consideradas pelo modelo. As regras verificam as condições necessárias para que ocorram alterações de estado durante a execução.

Nos mosquitos, a condição epidemiológica também influencia sua participação nos eventos de transmissão. A implementação deve distinguir a aquisição da infecção da condição em que o agente pode participar efetivamente da propagação representada pelo modelo.

Os eventos de transmissão são tratados como interações entre agentes humanos e vetoriais. A ocorrência de uma interação não implica necessariamente uma transmissão, pois as regras podem utilizar condições adicionais e componentes probabilísticos.

A implementação dessas regras deve manter correspondência com a especificação conceitual apresentada no capítulo anterior. Por esse motivo, as transições entre estados e suas condições constituem pontos importantes dos procedimentos de verificação.

\section{Implementação da dinâmica vetorial}

A dinâmica vetorial envolve os comportamentos dos mosquitos adultos e os processos relacionados às formas imaturas e aos criadouros.

Os mosquitos adultos possuem regras relacionadas ao deslocamento, à alimentação e à interação com a população humana. Essas regras determinam parte das oportunidades de transmissão representadas durante a simulação.

Os processos anteriores à vida adulta são organizados separadamente. A utilização de componentes destinados às formas imaturas permite representar etapas do ciclo vetorial sem concentrar todos os processos em uma única espécie de agente.

Os criadouros constituem a ligação entre determinados processos ambientais e a dinâmica da população vetorial. Suas condições podem ser modificadas de acordo com as regras ambientais utilizadas pelo modelo.

A implementação representa uma abstração do ciclo biológico real. Dessa forma, os comportamentos programados devem ser interpretados de acordo com as regras e parâmetros efetivamente definidos no simulador, sem assumir que todos os processos biológicos são representados individualmente.

\section{Controle temporal e climático}

O simulador utiliza um mecanismo global para controlar a evolução do tempo e organizar a execução dos diferentes processos.

A escala temporal da simulação é composta por ciclos computacionais. Diferentes comportamentos podem ser executados em frequências distintas, conforme as condições definidas nas regras do modelo.

O controle temporal é utilizado para coordenar processos como movimentação, evolução dos estados, atualização das condições ambientais e registro das informações produzidas.

Os dados climáticos são incorporados como entradas utilizadas pelas regras ambientais. As informações disponíveis podem ser processadas de acordo com o período representado e aplicadas aos componentes do modelo que dependem das condições de temperatura ou precipitação.

A utilização de dados climáticos no simulador não significa que cada variável ambiental seja aplicada diretamente a todos os agentes. Os efeitos dependem das regras implementadas e dos componentes aos quais essas informações são associadas.

A correspondência entre os ciclos computacionais e o tempo representado deve ser mantida de forma consistente, especialmente porque os resultados e os indicadores precisam ser interpretados de acordo com sua escala temporal.

\section{Interface e configuração dos parâmetros}

A plataforma GAMA permite disponibilizar parâmetros para a configuração das execuções. Esses parâmetros possibilitam alterar determinadas condições do modelo sem modificar diretamente os arquivos-fonte.

A interface do experimento é utilizada para configurar valores relacionados às condições da simulação e acompanhar indicadores durante sua execução.

A separação entre parâmetros configuráveis e regras implementadas contribui para a realização de experimentos controlados. Entretanto, para que um parâmetro seja efetivamente utilizado em um experimento, é necessário verificar sua ligação com os processos correspondentes da simulação.

Os valores utilizados em cada execução devem ser registrados como parte da configuração experimental. Esse registro é importante para permitir a reprodução e a comparação entre diferentes cenários.

\section{Registro e exportação dos resultados}

Durante a execução, o simulador registra indicadores relacionados ao comportamento das populações e às condições da simulação.

Os dados produzidos podem incluir informações sobre os estados epidemiológicos das populações humanas e vetoriais, além de outros indicadores definidos pelas regras do modelo.

O projeto possui mecanismos para a persistência das informações produzidas em arquivos de saída. A utilização de arquivos estruturados permite que os resultados sejam posteriormente analisados por outras ferramentas.

Para que os resultados sejam rastreáveis, é necessário associar os dados exportados às condições de cada execução. Informações como configuração utilizada, identificação da repetição, semente e versão do código contribuem para a interpretação e reprodução dos experimentos.

A existência de uma rotina de exportação não garante, isoladamente, a consistência dos indicadores registrados. As relações entre os totais das populações e as contagens por estado devem ser verificadas antes que as saídas sejam utilizadas para conclusões sobre o comportamento epidemiológico do modelo.

\section{Rastreabilidade entre modelo e implementação}

A rastreabilidade estabelece uma relação entre os elementos definidos no modelo conceitual e os componentes utilizados em sua implementação.

No simulador, essa relação pode ser organizada a partir das entidades, regras e processos apresentados anteriormente. Os agentes humanos correspondem aos componentes responsáveis pela representação da população sintética, enquanto os agentes vetoriais implementam os processos relacionados aos mosquitos.

Os componentes ambientais e territoriais correspondem às estruturas utilizadas para representar o espaço e as condições associadas aos criadouros. Os processos globais, por sua vez, concentram elementos compartilhados, como inicialização, controle temporal e registro de informações.

A rastreabilidade também permite identificar situações em que uma regra conceitual precisa ser confirmada por meio da análise ou da execução do código. Essa distinção é importante para evitar que um comportamento planejado seja apresentado como uma funcionalidade comprovadamente executada.

A documentação das relações entre o modelo conceitual e seus componentes computacionais facilita futuras alterações e contribui para os procedimentos de verificação.


\chapter{RESULTADOS E AVALIAÇÃO}

Este capítulo apresenta os resultados obtidos durante a análise do simulador e avalia as evidências disponíveis sobre seu funcionamento. A apresentação distingue três situações: resultados efetivamente registrados, verificações realizadas sobre esses registros e procedimentos experimentais que ainda necessitam de execução ou documentação adicional.

Essa distinção é necessária porque a existência de séries exportadas pelo simulador demonstra que determinadas rotinas produziram dados, mas não comprova, isoladamente, a correção de todas as regras implementadas ou a validade epidemiológica do modelo.

As análises apresentadas neste capítulo devem ser interpretadas de acordo com a versão do código e dos arquivos de saída efetivamente utilizados. Alterações posteriores na implementação podem modificar os resultados produzidos pelo simulador.

\section{Resultados da verificação}

A verificação inicial concentrou-se na análise da consistência dos indicadores registrados pelo simulador e na relação entre os valores exportados e as populações declaradas.

Um dos registros disponíveis apresenta uma população humana total de 1.191 agentes. Nesse mesmo registro, as quantidades de indivíduos nos estados suscetível, exposto, infectado e recuperado totalizam 1.191, indicando consistência aritmética entre esses indicadores.

Entretanto, foi identificada uma divergência nas contagens relacionadas à população vetorial. No registro analisado, o total de mosquitos é informado como 100, enquanto a soma dos agentes classificados nos estados suscetível, exposto e infectado resulta em 78.

Essa diferença demonstra a necessidade de revisar a lógica utilizada para contabilizar os estados epidemiológicos dos mosquitos e de verificar se todos os agentes presentes na população são classificados por algum dos indicadores exportados.

A identificação dessa divergência constitui um resultado de verificação da implementação. Ela não permite concluir, por si só, que a dinâmica vetorial esteja incorreta, mas evidencia que os indicadores disponíveis não devem ser utilizados como evidência de validade epidemiológica antes da investigação da causa da diferença.

Também foram identificadas limitações relacionadas à rastreabilidade das execuções. O arquivo de resultados analisado não contém, para cada registro, informações suficientes para identificar de forma inequívoca o experimento, a repetição, a semente utilizada, a versão do código e a configuração completa da execução.

Essas informações são importantes porque diferentes configurações podem produzir séries semelhantes ou distintas, e a ausência de metadados dificulta a reprodução dos resultados.

\section{Resultados do cenário de referência}

Os registros disponíveis demonstram que o simulador foi capaz de produzir séries temporais contendo informações sobre as populações humanas e vetoriais, além de variáveis ambientais.

Em um dos segmentos identificados no arquivo de resultados, a simulação inicia com 1.191 agentes humanos e 100 mosquitos. Ao longo dos registros, observa-se a alteração das contagens associadas aos estados epidemiológicos e à população vetorial.

Essas observações demonstram que a execução produziu comportamento dinâmico e que os mecanismos de registro foram capazes de gerar séries temporais. Entretanto, os dados disponíveis não possuem informações suficientes para associar de forma inequívoca cada segmento da série a uma configuração experimental específica.

Por esse motivo, os resultados existentes não são apresentados como um cenário de referência epidemiologicamente validado. Eles representam evidências exploratórias do funcionamento da rotina de execução e exportação do simulador.

Para a versão final da avaliação, o cenário de referência deverá ser definido por uma configuração completamente documentada. Essa configuração deve registrar, entre outros elementos, a versão do código, a versão da plataforma, os parâmetros utilizados, as condições iniciais, a duração da simulação e os dados de entrada.

\section{Comparação entre cenários}

A comparação entre cenários requer que cada execução possua uma configuração identificável e que as alterações realizadas sejam explicitamente registradas.

O código analisado apresenta configurações destinadas à execução de diferentes experimentos, incluindo variações de parâmetros relacionados à vacinação, à população vetorial e ao tamanho da população humana. Entretanto, a existência dessas configurações no código não demonstra, isoladamente, que todos os cenários foram executados ou que as alterações previstas foram efetivamente aplicadas durante as simulações.

Além disso, a análise estática identificou situações que exigem verificação experimental. Em uma das configurações, por exemplo, a existência de um parâmetro destinado à quantidade inicial de mosquitos não garante, por si só, que esse valor seja utilizado pela rotina responsável pela criação da população vetorial.

Dessa forma, os resultados atualmente disponíveis não permitem realizar uma comparação conclusiva entre os cenários planejados. A análise comparativa definitiva deve ser realizada após a execução controlada dos experimentos, com registro das configurações efetivamente utilizadas.

Para cada cenário, recomenda-se apresentar:

\begin{itemize}
    \item os parâmetros modificados;
    \item os valores utilizados;
    \item as condições mantidas constantes;
    \item o número de replicações;
    \item as sementes utilizadas;
    \item os indicadores selecionados para comparação.
\end{itemize}

A comparação deverá considerar a variabilidade entre as execuções e não apenas diferenças observadas em uma única trajetória.

\section{Variabilidade entre replicações}

A presença de componentes estocásticos faz com que diferentes execuções do simulador possam produzir trajetórias distintas, mesmo quando utilizam os mesmos parâmetros.

A avaliação da variabilidade requer múltiplas replicações de uma mesma configuração experimental. Para cada conjunto de execuções, podem ser analisadas medidas como média, desvio-padrão, quantis e outras estatísticas compatíveis com os indicadores produzidos.

Os registros atualmente disponíveis não possuem metadados suficientes para identificar de forma confiável todas as replicações e sementes associadas às séries analisadas. Consequentemente, não é possível utilizar essas saídas, isoladamente, para estimar de maneira conclusiva a variabilidade do modelo.

Uma campanha experimental futura deve utilizar uma configuração fixa e executar múltiplas replicações com sementes identificadas. A análise poderá então avaliar a dispersão dos resultados e verificar se as diferenças entre cenários são maiores que a variabilidade associada aos componentes aleatórios.

Também é necessário distinguir a variabilidade produzida pelo simulador da incerteza relacionada aos dados ou parâmetros utilizados. A dispersão entre execuções representa uma característica das condições experimentais e dos mecanismos estocásticos do modelo, não correspondendo automaticamente à incerteza sobre os casos reais de dengue.

\section{Avaliação computacional}

A avaliação computacional busca analisar características relacionadas ao custo de execução do simulador. Entre os indicadores de interesse podem estar o tempo necessário para concluir uma simulação, o consumo de memória e o comportamento do sistema diante do aumento da quantidade de agentes.

O código analisado possui uma configuração experimental relacionada à variação do tamanho da população. Entretanto, a presença dessa configuração não constitui, por si só, uma avaliação de desempenho.

Uma análise de escalabilidade requer medições realizadas sob condições controladas. Além do tamanho das populações, devem ser considerados fatores como a versão do software, a configuração da máquina utilizada, os dados de entrada e as condições de execução.

Também é necessário realizar mais de uma execução para cada configuração quando houver variabilidade relevante no tempo de processamento. Uma única medição não permite caracterizar adequadamente a dispersão do desempenho.

Portanto, a avaliação computacional definitiva depende da realização de medições sistemáticas. Caso esses experimentos sejam concluídos antes da versão final do trabalho, os resultados deverão ser apresentados nesta seção juntamente com as configurações utilizadas e os procedimentos de medição.

\section{Discussão das limitações}

Os resultados obtidos devem ser interpretados considerando as limitações identificadas durante a análise da implementação e das saídas disponíveis.

A primeira limitação está relacionada à consistência dos indicadores. A divergência observada entre o total de mosquitos e a soma dos estados vetoriais demonstra que determinadas contagens precisam ser verificadas antes que os resultados sejam utilizados para análises epidemiológicas.

A segunda limitação está relacionada à rastreabilidade dos experimentos. A ausência de informações como identificação da execução, semente efetiva e configuração completa impede a reprodução precisa das séries disponíveis e dificulta a comparação entre diferentes resultados.

Também existem limitações decorrentes das simplificações adotadas pelo modelo. A população humana é sintética, a representação espacial possui um nível específico de agregação e diversos processos biológicos e epidemiológicos são representados de forma simplificada.

Além disso, os resultados do simulador não devem ser interpretados automaticamente como previsões da ocorrência de dengue no município de São Paulo. Uma avaliação dessa natureza exigiria procedimentos adicionais de calibração e comparação com dados observados compatíveis.

A análise dos registros existentes demonstra que o simulador possui mecanismos capazes de gerar séries temporais e registrar indicadores durante sua execução. Entretanto, as evidências atualmente disponíveis não são suficientes para afirmar a validação epidemiológica do modelo ou sua capacidade preditiva.

As principais necessidades identificadas para o aprimoramento da avaliação são:

\begin{itemize}
    \item corrigir ou explicar as divergências entre os totais e os estados registrados;
    \item registrar metadados completos para cada execução;
    \item verificar o controle das sementes e a reprodutibilidade das execuções;
    \item realizar replicações controladas para avaliar a variabilidade;
    \item executar comparações entre cenários com parâmetros efetivamente aplicados;
    \item realizar medições sistemáticas de desempenho, quando essa avaliação fizer parte dos objetivos;
    \item estabelecer procedimentos específicos caso seja realizada comparação com dados epidemiológicos observados.
\end{itemize}






% Inserir os capítulos do trabalho nesta seção.

% ========================================================================
% PARTE PÓS-TEXTUAL
% ========================================================================

\postextual

% ========================================================================
% REFERÊNCIAS
% ========================================================================

\bibliography{bibliografia.bib}

% ========================================================================
% ANEXOS
% ========================================================================

\begin{anexosenv}

\chapter{Repositório do Projeto no GitHub}

O código-fonte desenvolvido para este Trabalho de Conclusão de Curso está disponível no seguinte repositório:

\begin{itemize}
    \item \url{https://github.com/Renatojcalpalhao/DengueSimu_Mesa}
\end{itemize}

\end{anexosenv}

\end{document}
